\documentclass[11pt,reqno]{amsart}
\usepackage[margin=1.15in]{geometry}
\usepackage{amsmath,amssymb,amsthm,mathtools}
\usepackage{booktabs}
\usepackage[colorlinks=true,linkcolor=blue,citecolor=blue,urlcolor=blue]{hyperref}
\usepackage{microtype}

\theoremstyle{plain}
\newtheorem{theorem}{Theorem}
\newtheorem{proposition}[theorem]{Proposition}
\newtheorem{corollary}[theorem]{Corollary}
\newtheorem{lemma}[theorem]{Lemma}
\newtheorem{conjecture}[theorem]{Conjecture}
\theoremstyle{definition}
\newtheorem{definition}[theorem]{Definition}
\newtheorem{remark}[theorem]{Remark}
\newtheorem{observation}[theorem]{Observation}

\DeclareMathOperator{\sqf}{sqf}
\newcommand{\Art}{\mathrm{A}}
\newcommand{\Z}{\mathbb{Z}}
\newcommand{\Q}{\mathbb{Q}}

\begin{document}

\title[Cross-base correlations of Artin primes]
{Cross-base correlations of Artin primes:\\ entanglement, exclusion, and the
structure of $p-1$}

\author{Josh Bald}
\address{Ontario, Canada}
\email{jpbald93@gmail.com}

\subjclass[2020]{Primary 11A07; Secondary 11N05, 11N13, 11Y16, 11Y60}
\keywords{Artin's conjecture, primitive roots, simultaneous primitive roots,
entanglement, quadratic characters, Kummer theory, computational number theory}

\date{\today}

\begin{abstract}
Say a prime $p$ is \emph{Artin base $a$} if $a$ is a primitive root modulo
$p$. This paper studies the joint behaviour of Artin status \emph{across
bases} at the same prime. We first record a deterministic exclusion law: if
$\sqf(ab)$ is the squarefree part of $ab$ and $c$ is any base with
$\sqf(c) = \sqf(ab)$, then no prime is simultaneously Artin for $a$, $b$, and
$c$ --- for instance, no prime has $2$, $5$, and $10$ all as primitive roots.
The proof is three lines from the multiplicativity of quadratic characters.
We then measure, for all $66$ pairs from the twelve bases
$\{2,3,5,6,7,10,11,13,15,17,21,29\}$ and all $50{,}847{,}532$ primes up to
$10^9$, the same-prime correlation $\phi(a,b)$ between the indicator variables
of Artin status. Three empirical findings emerge. First, the correlation is
positive for $64$ of the $66$ pairs (mean $\phi = 0.0352$, individual
$|z|$ up to $476$); the only negative pairs, $(2,10)$ and $(3,15)$, are
precisely pairs whose product is (a square times) another base --- the
entangled configurations. Second, conditioning on the number of prime factors
$\omega(p-1)$ removes only $55\%$ of the mean correlation, but conditioning on
a finer signature of $p - 1$ --- the $2$-adic valuation, the divisibility of
$p-1$ by each prime up to $13$, and the number of larger factors --- removes
$95\%$: the cross-base correlation is, almost entirely, a reflection of the
fact that all bases see the same factorisation of $p - 1$. Third, the residual
that survives fine conditioning is concentrated exactly on the multiplicatively
dependent pairs (mean absolute residual $0.018$ within dependent pairs,
versus $0.001$ for multiplicatively independent pairs), and is removed
completely by additionally conditioning on the quadratic-residue statuses of
the two bases --- identifying it as character-level entanglement, the same
mechanism as the deterministic exclusion law. Finally we evaluate the joint
densities predicted by the Matthews--Lenstra Kummer-degree heuristic for all
$66$ pairs: with no fitted parameters they match the measured densities to a
mean relative error of $0.037\%$, and they reproduce the sign reversal on the
entangled pairs. The deterministic exclusion law and the statistical residual
are thereby identified as two faces of one object, the degree collapse
contributed by $\sqf(ab)$.
\end{abstract}

\maketitle

\section{Introduction}
\label{sec:intro}

Artin's conjecture on primitive roots concerns one base at a time: for
non-square $a \neq -1$, the set
\[
  \Art_a = \{\, p \text{ prime} : a \text{ is a primitive root mod } p \,\}
\]
should have density given by a rational multiple of Artin's constant, as
proved under GRH by Hooley~\cite{Hooley1967}. The joint distribution for
\emph{several} bases at the same prime was treated by
Matthews~\cite{Matthews1976}, who computed (under GRH) the density of primes
for which several prescribed integers are simultaneously primitive roots, and
by Lenstra~\cite{Lenstra1977}, whose Galois-theoretic framework subsumes such
questions; see Moree's survey~\cite{Moree2012}. The qualitative message of
that theory is that the events $p \in \Art_a$ and $p \in \Art_b$ are
\emph{not} independent, and that the deviation from independence is governed
by the interaction of the Kummer--Galois extensions attached to $a$ and $b$
--- in modern language, by \emph{entanglement} between them.

That framework remains actively developed.
J\"arviniemi--Perucca--Sgobba~\cite{JarviniemiPeruccaSgobba2025} give a unified
treatment of Artin-type problems for finitely generated groups in number
fields; Kimmel~\cite{Kimmel2024} treats number-field and matrix variants;
Klurman--Shparlinski--Ter\"av\"ainen~\cite{KlurmanShparlinskiTeravainen2025}
obtain results on average over the base; Sgobba~\cite{Sgobba2025} establishes
unconditional results in cases where GRH can be circumvented; and
Goldmakher--Martin--P\'eringuey~\cite{GoldmakherMartinPeringuey2025} propose
refined conjectures on the distribution of the multiplicative order, which
bear on the structure of $p-1$ that plays the central role in
Section~\ref{sec:decomposition} below. None of this work addresses the joint
statistics of Artin status across several bases at a fixed prime, which is our
subject; we take~\cite{Matthews1976, Lenstra1977} as the source of the density
predictions tested in Section~\ref{sec:model}. As those predictions are
conditional on GRH, so is every comparison we make with them; our measurements
themselves are of course unconditional.

This paper is the third in a series studying correlations of Artin status
computationally at the scale of $10^9$~\cite{BaldI, BaldII}. The first
paper~\cite{BaldI} measured the correlation between the Artin statuses (base $10$) of
\emph{consecutive} primes and decomposed it into residue channels; the second
proved a classification of deterministic exclusion laws for consecutive Artin
primes in arbitrary bases and identified the conductor of the quadratic
character $\chi_a$ as the parameter governing the statistical correlation.
Here we complete the picture along the remaining axis: \emph{fix the prime,
vary the base}.

Two questions organise the paper.

\emph{(Q1) What is deterministically forbidden?} For consecutive primes, the
answer~\cite{BaldII} was a set of gap classes on which the
quadratic character is forced to flip. For a single prime and several bases,
the answer is an exclusion law inside \emph{multiplicative triples}:

\begin{theorem}[Same-prime exclusion law; see Theorem~\ref{thm:triple}]
Let $a, b$ be non-square positive integers and let $c$ be any non-square
positive integer with $\sqf(c) = \sqf(ab)$. Then no prime $p \nmid 2abc$ is
simultaneously Artin base $a$, Artin base $b$, and Artin base $c$.
\end{theorem}

For example, no prime has all of $2, 5, 10$ as primitive roots; no prime has
all of $3, 5, 15$; and, since $\sqf$ is all that matters, no prime has all of
$2, 5, 40$. Like the gap exclusion law of~\cite{BaldII}, the statement is
elementary --- it needs only the multiplicativity of the quadratic character
and the fact that a quadratic residue is not a primitive root --- and, like
that law, we have not found it stated explicitly in the literature, although
it is certainly implicit in the density computations
of~\cite{Matthews1976, Lenstra1977}, where the entangled character relation
$\chi_a \chi_b = \chi_{ab}$ is exactly what produces degenerate factors.

\emph{(Q2) What does the joint distribution actually look like, and what
explains it?} Here the data speak. We compute the exact $2 \times 2$
contingency table of $\bigl(\mathbf{1}[p \in \Art_a],\,
\mathbf{1}[p \in \Art_b]\bigr)$ over all $50{,}847{,}532$ primes
$p \le 10^9$, for all $66$ pairs from the base set
\[
  \{2, 3, 5, 6, 7, 10, 11, 13, 15, 17, 21, 29\},
\]
chosen to contain many multiplicative triples ($15$ of the $66$ pairs lie in
a triple $\sqf(ab) = \sqf(c)$ within the set) alongside multiplicatively
independent pairs. The findings, developed in
Sections~\ref{sec:measurement}--\ref{sec:entanglement}:

\begin{enumerate}
  \item \emph{Positivity.} $\phi(a,b) > 0$ for $64$ of $66$ pairs, with mean
    $0.0352$ and $z$-scores up to $476$: knowing that $a$ is a primitive root
    mod $p$ makes it \emph{more} likely that $b$ is. The largest correlations
    all involve base $5$ (e.g.\ $\phi(5,13) = 0.0667$:
    $P(\Art_{13} \mid \Art_5) = 0.416$ against
    $P(\Art_{13} \mid \lnot\Art_5) = 0.350$).
  \item \emph{The exceptions are the entangled pairs.} The only two negative
    correlations are $\phi(2,10) = -0.0302$ and $\phi(3,15) = -0.0303$ ---
    precisely the pairs $(a, b)$ in our set with $b = a \cdot c$ for another
    base $c$ of the set (up to squares). Multiplicative dependence reverses
    the sign of the association.
  \item \emph{The positive correlation is $p-1$ structure, not entanglement.}
    Conditioning on $\omega(p-1)$ removes $55\%$ of the mean correlation;
    conditioning on a fine signature of $p-1$ --- namely
    $\bigl(\min(v_2(p-1), 4),\ \text{the set of primes } q \le 13 \text{
    dividing } p-1,\ \min(\omega_{>13}(p-1), 3)\bigr)$ --- removes $95\%$,
    leaving a mean residual of $0.0019$. All bases are tested against the
    same integer $p - 1$; a prime whose $p-1$ is ``primitive-root friendly''
    (few small factors, small $2$-part) is friendlier for \emph{every} base
    simultaneously. This common-cause structure, not any interaction between
    the bases, is the bulk of the correlation.
  \item \emph{What survives is entanglement, and it is exactly localised.}
    After fine conditioning, the residual correlation among multiplicatively
    \emph{dependent} pairs (mean $|\phi_{\mathrm{sig}}| \approx 0.005$, and
    mean absolute residual $0.0179$, with individual values reaching
    $-0.034$ for $(2,10)$ and $-0.026$ for $(5,15)$) is eighteen times that of
    independent pairs ($0.0010$); and additionally conditioning on the
    quadratic-residue statuses $(\chi_a(p), \chi_b(p))$ removes it entirely,
    flattening dependent and independent pairs to a common $+0.005$. The
    residual is thus identified, empirically, as the quadratic-character
    entanglement $\chi_a \chi_b = \chi_{\sqf(ab)}$ --- the statistical shadow
    of Theorem~\ref{thm:triple}.
  \item \emph{Theory reproduces all of it.} Evaluating the Matthews--Lenstra
    Kummer-degree heuristic~\eqref{eq:model} for each pair
    (Section~\ref{sec:model}) predicts the measured joint densities to a mean
    relative error of $0.037\%$ over the $66$ pairs, with no fitted
    parameters, including the negative $\phi$ of the two entangled pairs. The
    entangled degree collapse in the model is precisely the character relation
    behind Theorem~\ref{thm:triple}.
\end{enumerate}

The theme of the series persists: at every level, the joint behaviour of
Artin primes decomposes into a small deterministic core (here, the triple
exclusion law) plus a statistical component with an identifiable arithmetic
carrier (here, the shared factorisation of $p-1$, with entanglement as a
measurable second-order correction).

\section{The same-prime exclusion law}
\label{sec:theorem}

Throughout, for a non-square positive integer $a$ we write $d_a = \sqf(a)$
for its squarefree part and $\chi_a = \bigl(\tfrac{d_a}{\cdot}\bigr)$ for the
quadratic character attached to $\Q(\sqrt{a})$; thus $\chi_a(p) = +1$ exactly
when $a$ is a quadratic residue modulo $p$ (for $p \nmid 2a$), and
\begin{equation}
\label{eq:obstruction}
  p \in \Art_a \implies \chi_a(p) = -1,
\end{equation}
since a quadratic residue generates a subgroup of index at least $2$ in
$(\Z/p\Z)^\times$.

\begin{theorem}[Triple exclusion]
\label{thm:triple}
Let $a, b, c$ be non-square positive integers with
$\sqf(c) = \sqf(ab)$, and let $p \nmid 2abc$ be prime. Then $p$ is not
simultaneously Artin base $a$, Artin base $b$, and Artin base $c$.
\end{theorem}

\begin{proof}
For any prime $p \nmid 2ab$ one has, by complete multiplicativity of the
Legendre symbol in the numerator,
$\chi_a(p)\,\chi_b(p) = \bigl(\tfrac{d_a d_b}{p}\bigr) =
\bigl(\tfrac{\sqf(ab)}{p}\bigr) = \chi_c(p)$,
since $d_a d_b$ and $\sqf(ab)$ differ by a square. If $p$ were Artin for all
three bases, then by~\eqref{eq:obstruction} all three character values would
be $-1$, giving $(-1)(-1) = -1$, a contradiction.
\end{proof}

\begin{corollary}
\label{cor:twoofthree}
Within any multiplicative triple --- for instance $(2, 5, 10)$, $(3, 5, 15)$,
$(3, 7, 21)$, $(2, 3, 6)$ --- a prime can be Artin for at most two of the
three bases.
\end{corollary}

\begin{remark}
The exclusion is sharp: primes Artin for exactly two of $(2,5,10)$ exist in
abundance. The smallest is $p = 7$, where $5$ and $10$ are both primitive
roots (each of order $6$) while $2$ has order $3$; likewise $p = 23$, where
$5$ and $10$ are primitive roots and $2$ is a quadratic residue of order
$11$. More systematically, in our data the pair
$(2, 10)$ has $P(\Art_{10} \mid \Art_2) = 0.355$: two-of-three is common,
three-of-three is impossible.
\end{remark}

\begin{remark}
Theorem~\ref{thm:triple} is the same-prime analogue of the gap exclusion law
for consecutive primes proved in~\cite{BaldII}: both are
instances of a parity obstruction in the group generated by the quadratic
characters involved. There the two characters were $\chi_a$ evaluated at $p$
and at $p + g$, and the obstruction was arranged by the shift $g$; here the
characters belong to different bases at the same prime, and the obstruction
is arranged by multiplicative dependence. The Galois-theoretic density
computations of Matthews~\cite{Matthews1976} and Lenstra~\cite{Lenstra1977}
encode this relation implicitly (the degree of the compositum of the relevant
Kummer extensions drops); the pointwise statement seems worth recording.
\end{remark}

\section{Data and computation}
\label{sec:data}

By a segmented sieve we enumerated all primes $p \le 10^9$
($50{,}847{,}534$ primes, of which $50{,}847{,}532$ satisfy $p \nmid 2\cdot
\text{(base)}$ for all twelve bases and enter the analysis). For each prime we
factored $p - 1$ once by trial division and computed the Artin indicator for
each of the twelve bases by the standard criterion
($a^{(p-1)/\ell} \not\equiv 1 \bmod p$ for every prime $\ell \mid p-1$, in
$128$-bit arithmetic). One streaming pass accumulates, for each of the $66$
unordered base pairs, the $2 \times 2$ contingency table of joint Artin
status --- globally, and stratified by (i) $\omega(p-1)$ and (ii) the fine
signature
\[
  \mathrm{sig}(p) = \Bigl(\min(v_2(p-1), 4),\;\;
  \{\, q \le 13 : q \mid p-1 \,\},\;\; \min(\omega_{>13}(p-1), 3)\Bigr),
\]
a partition of the primes into at most $5 \cdot 32 \cdot 4 = 640$ cells. The
signature records exactly the information about $p-1$ that the Artin
criterion consults for small bases: the $2$-part and the small odd prime
divisors. A separate run, also to $10^9$ ($50{,}847{,}524$ primes), stratified
each pair additionally by the quadratic-residue statuses
$\bigl(\chi_a(p), \chi_b(p)\bigr)$ of the two bases on top of the signature.
The whole computation runs in under two hours on a single core of a
commodity machine ($2$-core Intel Xeon E5-2673~v4, $8$~GB RAM); all counts
are exact.

As validation, the marginal Artin densities matched the Hooley densities to
four decimal places for every base --- e.g.\ $0.373993$ for the bases with
squarefree part $\not\equiv 1 \pmod 4$, against Artin's constant
$0.373956$, and $0.393642$ for base $5$, whose density carries the classical
correction factor $C \cdot \tfrac{20}{19}$ for
$d \equiv 1 \pmod 4$ (see~\cite{Hooley1967, Moree2012}).

\section{The cross-base correlation}
\label{sec:measurement}

For each pair write $\phi(a,b)$ for the Pearson correlation of the indicator
variables (the $\phi$ coefficient of the $2\times2$ table). At
$N \approx 5.08 \times 10^7$ primes, the standard error of $\phi$ under
independence is $N^{-1/2} \approx 0.00014$; as in the previous papers of the
series, the enormous resulting $z$-scores certify only that sampling noise is
negligible, and the relevant quantity is the effect size.

\begin{table}[t]
\centering
\caption{Same-prime cross-base Artin correlation at $x = 10^9$: summary and
extremes. ``Dependent'' means the pair lies in a multiplicative triple within
the base set ($\sqf(ab) = \sqf(c)$).}
\label{tab:summary}
\begin{tabular}{lrr}
\toprule
 & pairs & mean $\phi$ \\
\midrule
all pairs & $66$ & $+0.0352$ \\
multiplicatively independent & $51$ & $+0.0363$ \\
multiplicatively dependent & $15$ & $+0.0314$ \\
\midrule
largest: $(5,13)$ & & $+0.0667$ \\
smallest: $(3,15)$ & & $-0.0303$ \\
second smallest: $(2,10)$ & & $-0.0302$ \\
\bottomrule
\end{tabular}
\end{table}

\begin{table}[t]
\centering
\caption{The ten largest $|\phi|$ at $x = 10^9$, with the values after
conditioning on $\omega(p-1)$ and on the fine signature $\mathrm{sig}(p)$
(pair-weighted mean of within-cell $\phi$; cells with fewer than $200$
primes excluded). ``Dep.''\ marks multiplicatively dependent pairs.}
\label{tab:top}
\begin{tabular}{rrrrrc}
\toprule
$a$ & $b$ & $\phi$ & $\phi \mid \omega$ & $\phi \mid \mathrm{sig}$ & Dep. \\
\midrule
$5$ & $13$ & $+0.0667$ & $+0.0200$ & $+0.0003$ & \\
$5$ & $17$ & $+0.0649$ & $+0.0218$ & $+0.0025$ & \\
$5$ & $29$ & $+0.0631$ & $+0.0246$ & $+0.0007$ & \\
$5$ & $10$ & $+0.0624$ & $+0.0285$ & $-0.0125$ & \checkmark \\
$5$ & $6$  & $+0.0623$ & $+0.0292$ & $+0.0005$ & \\
$5$ & $11$ & $+0.0623$ & $+0.0300$ & $+0.0004$ & \\
$5$ & $15$ & $+0.0622$ & $+0.0297$ & $-0.0255$ & \checkmark \\
$3$ & $5$  & $+0.0622$ & $+0.0284$ & $+0.0015$ & \checkmark \\
$5$ & $7$  & $+0.0622$ & $+0.0297$ & $+0.0002$ & \\
$2$ & $5$  & $+0.0621$ & $+0.0336$ & $+0.0004$ & \checkmark \\
\bottomrule
\end{tabular}
\end{table}

Three observations, in increasing order of depth.

\begin{observation}[Positivity]
\label{obs:positive}
$\phi(a, b) > 0$ for $64$ of the $66$ pairs. Artin statuses of different
bases at the same prime \emph{attract}.
\end{observation}

The sign is the opposite of the consecutive-prime correlation measured in the
first two papers of the series~\cite{BaldI, BaldII} (where $\delta < 0$ uniformly), and the
mechanism is transparent: there the two primes had \emph{different} (and
LOS-anticorrelated) values of $p - 1$; here every base is tested against the
\emph{same} $p - 1$. If $p - 1$ has few small prime factors and a small
$2$-part, the primitive-root condition is easier for every base at once.
Formally, the Artin indicators share the common cause $\mathrm{sig}(p)$.

\begin{observation}[The negative pairs are the entangled ones]
\label{obs:negative}
The only pairs with $\phi < 0$ are $(2, 10)$ and $(3, 15)$: in each, the
product of the two bases is, up to a square, the third base of a
multiplicative triple contained in the set ($2 \cdot 10 = 4 \cdot 5$,
$3 \cdot 15 = 9 \cdot 5$). For $(2, 10)$:
$P(\Art_{10} \mid \Art_2) = 0.3551$ versus
$P(\Art_{10} \mid \lnot\Art_2) = 0.3853$.
\end{observation}

The sign reversal has a character-theoretic origin. For a dependent pair,
$\chi_a \chi_b = \chi_5$ (in both our examples), so on the half of primes
with $\chi_5(p) = +1$ the values $\chi_a(p)$ and $\chi_b(p)$ are forced
\emph{equal} --- fine for joint Artin status --- but on the half with
$\chi_5(p) = -1$ they are forced \emph{opposite}, and then at most one of the
two bases can be Artin: a pointwise exclusion on half the primes, diluted
into the observed negative correlation. This is Theorem~\ref{thm:triple}
operating statistically through the third character.

\section{Decomposition: the correlation is the structure of \texorpdfstring{$p-1$}{p-1}}
\label{sec:decomposition}

\begin{table}[t]
\centering
\caption{Mean cross-base correlation after conditioning, all $66$ pairs at
$x = 10^9$.}
\label{tab:decomp}
\begin{tabular}{lrr}
\toprule
conditioning & mean $\phi$ & removed \\
\midrule
none & $0.0352$ & --- \\
$\omega(p-1)$ & $0.0158$ & $55\%$ \\
$\mathrm{sig}(p)$: $\bigl(v_2,\ \{q \le 13\},\ \omega_{>13}\bigr)$
  & $0.0019$ & $95\%$ \\
\bottomrule
\end{tabular}
\end{table}

Conditioning on $\omega(p-1)$ alone --- the crude size of the factorisation
--- removes barely half of the correlation. This is worth emphasising because
$\omega$ is the statistic through which the consecutive-prime papers routed
the mechanism, and it is genuinely insufficient here: two primes with the same
$\omega(p-1)$ can present very different Artin conditions (it matters
\emph{which} primes divide $p-1$, and with what $2$-adic valuation, not how
many). The fine signature --- $2$-adic valuation capped at $4$, the exact set
of primes $\le 13$ dividing $p - 1$, and the count of larger prime factors
capped at $3$ --- is the information the Artin criterion actually consults
for small bases, and conditioning on it removes $95\%$ of the mean
correlation (Table~\ref{tab:decomp}). That the fine structure of $p-1$, rather
than its coarse size, is what governs Artin behaviour is consistent with the
refined order-distribution conjectures
of~\cite{GoldmakherMartinPeringuey2025}, where the same distinction between
$\omega(p-1)$ and the actual divisor structure is what refines Artin's original
heuristic.

We regard this as the main statistical finding: \emph{cross-base Artin
correlation at a single prime is, to first order, not an interaction between
the bases at all}. It is a common-cause effect of the shared integer $p - 1$.
Any model that treats $\Art_a$ and $\Art_b$ as conditionally independent
given the small-prime structure of $p-1$ reproduces $95\%$ of the observed
association.

\section{The residual is entanglement}
\label{sec:entanglement}

The $5\%$ that survives is not noise, and it is not uniformly distributed:
it sits precisely on the multiplicatively dependent pairs.

\begin{table}[t]
\centering
\caption{Residual correlation after conditioning on $\mathrm{sig}(p)$ and
after additionally conditioning on the pair of quadratic-residue statuses
$(\chi_a(p), \chi_b(p))$, split by multiplicative dependence. All values at
$x = 10^9$ over $50{,}847{,}524$ primes.}
\label{tab:residual}
\begin{tabular}{lrrr}
\toprule
group & mean $\phi$ & mean $\phi \mid \mathrm{sig}$
      & mean $\phi \mid \mathrm{sig}, \mathrm{QR}$ \\
\midrule
multiplicatively dependent ($15$ pairs)   & $+0.0314$ & $+0.0049$ & $+0.0049$ \\
\quad mean \emph{absolute} residual        &           & $\;\;0.0179$ & \\
multiplicatively independent ($51$ pairs) & $+0.0363$ & $+0.0010$ & $+0.0052$ \\
\quad mean \emph{absolute} residual        &           & $\;\;0.0010$ & \\
\midrule
all $66$ pairs                            & $+0.0352$ & $+0.0019$ & $+0.0052$ \\
\bottomrule
\end{tabular}
\end{table}

Two facts pin down the mechanism, both measured over the full $10^9$ range.

First, after fine conditioning the dependent pairs retain substantial
residuals of the sign predicted by the character relation of
Observation~\ref{obs:negative}: their mean \emph{absolute} residual is
$0.0179$, eighteen times the $0.0010$ of the $51$ independent pairs. Indeed
the six most negative residuals in the entire collection are \emph{exactly}
six dependent pairs and nothing else:
\[
  (2,10)\!:\, -0.0343, \quad (5,15)\!:\, -0.0255, \quad (6,10)\!:\, -0.0144,
  \quad (5,10)\!:\, -0.0125, \quad (7,21)\!:\, -0.0060, \quad (10,15)\!:\, -0.0049 .
\]
Note that $(2,10)$ and $(5,15)$ are more negative \emph{after} conditioning
than before, and that $(5,10)$, $(5,15)$, $(6,10)$, $(10,15)$ and $(7,21)$
change sign under conditioning: the common-cause channel had been masking an
underlying repulsion.

Second, additionally conditioning on the pair of quadratic-residue statuses
$(\chi_a(p), \chi_b(p))$ removes this residual completely. Every one of the
six negative values above becomes positive and small
($+0.0031$, $+0.0066$, $+0.0045$, $+0.0065$, $+0.0040$, $+0.0048$
respectively), and the dependent and independent groups become
indistinguishable at $+0.0049$ and $+0.0052$ (Table~\ref{tab:residual}).
That common small positive floor of about $+0.005$ is present for all $66$
pairs and is plausibly attributable to higher-order (cubic and quartic)
conditions shared through $p - 1$ beyond what our signature records.

Conditioning on the quadratic characters is precisely what removes
character-level entanglement; that it flattens exactly the pairs with
$\chi_a \chi_b = \chi_{\sqf(ab)}$, and only those, identifies the residual as
the statistical trace of the entangled relation --- the same object that makes
Theorem~\ref{thm:triple} true.

The picture that emerges is cleanly layered:
\[
  \underbrace{\text{common cause } p-1}_{\approx 95\%}
  \;+\;
  \underbrace{\text{quadratic entanglement } \chi_a\chi_b = \chi_{\sqf(ab)}}_{\text{residual, sign-definite, localised}}
  \;+\;
  \underbrace{\text{higher-order sharing}}_{\text{small positive floor}}.
\]

\section{Comparison with the predicted joint densities}
\label{sec:model}

The measurements of Sections~\ref{sec:measurement}--\ref{sec:entanglement}
identify the mechanisms empirically. We now check them against theory
quantitatively, by computing the joint densities predicted by the
Matthews--Lenstra framework and comparing with the observed counts pair by
pair.

Write $d_a = \sqf(a)$. The heuristic, in the form given
by~\cite{Matthews1976} and systematised in~\cite{Lenstra1977}, is an inclusion
--exclusion over the degrees of the Kummer extensions that detect the
conditions ``$a$ is an $\ell$-th power residue'':
\begin{equation}
\label{eq:model}
  \mathbf{P}\bigl(p \in \Art_a \cap \Art_b\bigr)
  \;=\; \sum_{m, n \text{ squarefree}}
  \frac{\mu(m)\,\mu(n)}
       {\bigl[\Q(\zeta_L, a^{1/m}, b^{1/n}) : \Q\bigr]},
  \qquad L = \mathrm{lcm}(m,n),
\end{equation}
and the entanglement enters through the degree in the denominator. Generically
$[\Q(\zeta_L, a^{1/m}, b^{1/n}):\Q] = \varphi(L)\,m\,n$, but the degree
\emph{drops by a factor of $2$ for each quadratic class in the group}
\[
  V(m,n) = \bigl\langle\, d_a \text{ (if $2 \mid m$)},\;
                          d_b \text{ (if $2 \mid n$)} \,\bigr\rangle
  \subseteq \Q^\times/(\Q^\times)^2
\]
whose quadratic field lies inside $\Q(\zeta_L)$, i.e.\ whose conductor divides
$L$. The crucial point for this paper is that when both $2 \mid m$ and
$2 \mid n$ the group $V$ contains not only $d_a$ and $d_b$ but their product
$d_a d_b \sim \sqf(ab)$: the entangled third character of
Theorem~\ref{thm:triple} appears in the model as an extra degree collapse,
exactly for the multiplicatively dependent pairs. We evaluate~\eqref{eq:model}
by summing the ramified part exactly over all squarefree $m, n$ supported on
$S = \{2\} \cup \{q : q \mid d_a d_b\}$ and folding the rest into the Euler
product $\prod_{q \notin S}\bigl(1 - \tfrac{2q-1}{q^2(q-1)}\bigr)$,
truncated at $3 \times 10^6$.

As a check on the implementation, the same code specialised to a single base
reproduces the classical Hooley densities: $0.373956$ for bases with
$d \not\equiv 1 \pmod 4$ (Artin's constant), $0.393638$ for $d = 5$,
$0.376368$ for $d = 13$, against measured $0.373993$, $0.393642$, $0.376369$.

\begin{table}[t]
\centering
\caption{Predicted versus measured joint Artin densities, selected pairs at
$x = 10^9$. ``Dep.''\ marks multiplicatively dependent pairs. The final two
rows are the two negatively correlated pairs.}
\label{tab:model}
\begin{tabular}{rrrrrrrc}
\toprule
$a$ & $b$ & measured $j$ & model $j$ & rel.\ err.
    & measured $\phi$ & model $\phi$ & Dep. \\
\midrule
$5$  & $13$ & $0.163939$ & $0.164671$ & $+0.45\%$ & $+0.0667$ & $+0.0698$ & \\
$5$  & $10$ & $0.161962$ & $0.161922$ & $-0.02\%$ & $+0.0624$ & $+0.0623$ & \checkmark \\
$5$  & $15$ & $0.161934$ & $0.161922$ & $-0.01\%$ & $+0.0622$ & $+0.0623$ & \checkmark \\
$3$  & $7$  & $0.149965$ & $0.149971$ & $+0.00\%$ & $+0.0432$ & $+0.0433$ & \checkmark \\
$13$ & $17$ & $0.150249$ & $0.150314$ & $+0.04\%$ & $+0.0383$ & $+0.0386$ & \\
$2$  & $3$  & $0.147361$ & $0.147349$ & $-0.01\%$ & $+0.0321$ & $+0.0321$ & \checkmark \\
$2$  & $6$  & $0.147334$ & $0.147349$ & $+0.01\%$ & $+0.0320$ & $+0.0321$ & \checkmark \\
$7$  & $21$ & $0.144754$ & $0.144728$ & $-0.02\%$ & $+0.0238$ & $+0.0238$ & \checkmark \\
$15$ & $21$ & $0.144725$ & $0.144728$ & $+0.00\%$ & $+0.0237$ & $+0.0238$ & \\
\midrule
$3$  & $15$ & $0.132769$ & $0.132776$ & $+0.01\%$ & $-0.0303$ & $-0.0302$ & \checkmark \\
$2$  & $10$ & $0.132776$ & $0.132776$ & $+0.00\%$ & $-0.0303$ & $-0.0302$ & \checkmark \\
\bottomrule
\end{tabular}
\end{table}

\begin{observation}[The model reproduces the measurements]
\label{obs:model}
Over all $66$ pairs at $x = 10^9$, the mean absolute relative error between
the predicted and measured joint densities is $0.037\%$, with maximum
$0.67\%$ (attained at $(5,21)$). In particular the model predicts the sign
reversal: for the entangled pairs $(2,10)$ and $(3,15)$ it gives
$j = 0.132776$ and $\phi = -0.0302$, against measured $j = 0.132776$,
$0.132769$ and $\phi = -0.0303$ for both.
\end{observation}

This is the quantitative counterpart of Sections~\ref{sec:measurement}
and~\ref{sec:entanglement}, and it settles the interpretation. The
$66$-pair pattern --- positive correlation almost everywhere, negative
precisely on the two entangled pairs, with magnitudes ordered as observed
--- is not merely \emph{consistent with} the Matthews--Lenstra framework; it is
reproduced by it to within a few parts in $10^4$, with no fitted parameters.
Correspondingly, the deterministic law of Theorem~\ref{thm:triple} and the
statistical residual isolated in Section~\ref{sec:entanglement} are two faces
of one object: the degree collapse contributed by $\sqf(ab)$ in $V(m,n)$.

\begin{remark}
The largest discrepancies ($0.67\%$ at $(5,21)$, $0.11\%$ at $(13,21)$)
involve base $21$, the only base in our set whose squarefree part has two odd
prime factors; the residual is consistent with the truncation of the Euler
tail at $3 \times 10^6$ interacting with the larger ramified set
$S = \{2,3,5,7\}$, and we have not pursued it. Ordinary Hooley-type errors of
this size are also expected at $x = 10^9$, since the convergence in Hooley's
theorem is only logarithmic.
\end{remark}

\section{Discussion}
\label{sec:discussion}

\subsection*{Relation to the first two papers of the series~\cite{BaldI, BaldII}}
The trilogy now covers both axes of the joint distribution of Artin status.
Along the prime axis (consecutive $p_n, p_{n+1}$, one base), correlation is
negative, driven by the Lemke Oliver--Soundararajan bias~\cite{LOS2016}
propagating through the two distinct integers $p_n - 1$ and $p_{n+1} - 1$,
with deterministic exclusion classes in the gap. Along the base axis (one
prime, two bases), correlation is positive, driven by the single shared
integer $p - 1$, with deterministic exclusion inside multiplicative triples.
In both cases the deterministic law is a parity statement about quadratic
characters, and in both cases it is statistically dominated by a
factorisation-of-$p-1$ mechanism --- diluted to invisibility in the global
average there, absorbed almost entirely by conditioning here.

\subsection*{Future work}
(1) A rigorous version of Section~\ref{sec:model}: the comparison there is
between measurement and a heuristic evaluated numerically, and the error terms
--- both the Hooley-type error at $x = 10^9$ and the truncation of our Euler
tail --- are not controlled. Making the agreement into a theorem (conditionally
on GRH, in the manner of~\cite{Matthews1976}) would require uniform error
bounds we do not attempt. (2) The mixed axis: $\phi$ between
$\Art_a(p_n)$ and $\Art_b(p_{n+1})$ for $a \neq b$ --- the data pipeline
extends immediately, and the prediction from the present findings is a small
negative correlation governed by the two conductors jointly. (3) Higher-order
entanglement: triples of bases conditioned on quadratic data should expose
the cubic layer ($\ell = 3$ Kummer conditions), with $\Q(\zeta_3)$ playing
the role $\Q(i)$ and $\Q(\sqrt{2})$ played here.

\section*{Acknowledgements}
The author thanks the developers of the open-source tools used in this work
(GCC, Python, \texttt{sympy}, \texttt{matplotlib}).

\subsection*{Use of generative AI}
The author used Anthropic's Claude (models Claude Opus 4.5 and Claude
Opus 5) as an assistant: for drafting and revising exposition, for writing
and debugging the sieve and analysis programs, and for exploratory
discussion. All computations were executed and verified on the author's own
hardware from source code the author has inspected; the proof of
Theorem~\ref{thm:triple} has been checked by hand; every reference has been
verified against the published record. The author is solely responsible for
the content.

\section*{Declaration of interest}
The author reports there are no competing interests to declare.

\section*{Funding}
This research received no external funding.

\section*{Data availability}
All code and results are publicly available at
\url{https://github.com/jpbald93/consecutive-artin}, including the
single-pass multi-base programs (\texttt{crossbase.c},
\texttt{crossbase\_fine.c}, \texttt{crossbase\_qr.c}), the exact contingency
tables underlying every value quoted here, and the analysis scripts. The full
computation reproduces in under two hours on a single core.

\begin{thebibliography}{99}

\bibitem{BaldI}
J.~Bald,
\emph{Correlations between primitive root statuses of consecutive primes},
preprint (2026); code and data at
\url{https://github.com/jpbald93/consecutive-artin}.

\bibitem{BaldII}
J.~Bald,
\emph{Exclusion laws for consecutive Artin primes in arbitrary bases},
preprint (2026); code and data at
\url{https://github.com/jpbald93/consecutive-artin}.

\bibitem{GoldmakherMartinPeringuey2025}
L.~Goldmakher, G.~Martin, and P.~P\'eringuey,
\emph{Refinements of Artin's primitive root conjecture},
preprint (2025), \texttt{arXiv:2502.19601}.

\bibitem{Hooley1967}
C.~Hooley,
\emph{On Artin's conjecture},
J. Reine Angew. Math. \textbf{225} (1967), 209--220.

\bibitem{JarviniemiPeruccaSgobba2025}
O.~J\"arviniemi, A.~Perucca, and P.~Sgobba,
\emph{Unified treatment of Artin-type problems II},
Res. Number Theory \textbf{11} (2025), no.~42, 19~pp.

\bibitem{Kimmel2024}
N.~Kimmel,
\emph{Artin's primitive root conjecture in number fields and for matrices},
Acta Arith. \textbf{216} (2024), no.~4, 329--348.

\bibitem{KlurmanShparlinskiTeravainen2025}
O.~Klurman, I.~E. Shparlinski, and J.~Ter\"av\"ainen,
\emph{On Artin's conjecture on average and short character sums},
Bull. Lond. Math. Soc. \textbf{57} (2025), 2429--2443.

\bibitem{Lenstra1977}
H.~W. Lenstra, Jr.,
\emph{On Artin's conjecture and Euclid's algorithm in global fields},
Invent. Math. \textbf{42} (1977), 201--224.

\bibitem{LOS2016}
R.~J. Lemke~Oliver and K.~Soundararajan,
\emph{Unexpected biases in the distribution of consecutive primes},
Proc. Natl. Acad. Sci. USA \textbf{113} (2016), no.~31, E4446--E4454.

\bibitem{Matthews1976}
K.~R. Matthews,
\emph{A generalisation of Artin's conjecture for primitive roots},
Acta Arith. \textbf{29} (1976), no.~2, 113--146.

\bibitem{Moree2012}
P.~Moree,
\emph{Artin's primitive root conjecture --- a survey},
Integers \textbf{12A} (2012), Article A13, 100~pp.

\bibitem{Sgobba2025}
P.~Sgobba,
\emph{Unconditional results for Artin-type problems over number fields},
preprint (2025), \texttt{arXiv:2508.08996}.

\end{thebibliography}

\end{document}
